% The reaction admissibility graph for the nine families. Drawn here rather
% than lifted from the mean-field treatment: that graph is a size-resolved
% ladder over four families, and this model carries moments over nine, with a
% basal chain and a grain-boundary loop channel the ladder has no vertex or
% edge for.
%
% EVERY EDGE IS CHECKED AGAINST THE SOLVER, not against the prose. The earlier
% version of this figure omitted the cascade sources of C_2i, C_3i and of the
% prismatic vacancy loops; the mobile clustering ladder and its dissociation
% edges entirely; the homogeneous loop-nucleation channel that 2i+2i and i+3i
% feed; and the capture edges onto the basal family. Each of those appears in
% `rate_equations_core.h`, and each is now drawn.
%
% The four mobile species are boxed together so that capture onto a loop can be
% ONE arrow per family rather than four: the growth law sums over all four
% species with a sign set by whether the polarity matches the family's, so the
% bundle is faithful and the alternative is twelve crossing arrows.
\begin{figure}[htbp]
\centering
\begin{tikzpicture}[
  font=\small, >=Latex,
  mob/.style={draw, rounded corners=2pt, minimum height=7mm,
              minimum width=10mm, fill=blue!8, draw=blue!55},
  ipop/.style={draw, rounded corners=2pt, align=center, fill=red!8,
               draw=red!60, minimum height=8mm, inner sep=3pt},
  vpop/.style={draw, rounded corners=2pt, align=center, fill=cyan!10,
               draw=cyan!55!black, minimum height=8mm, inner sep=3pt},
  bpop/.style={draw, rounded corners=2pt, align=center, fill=violet!10,
               draw=violet!60, minimum height=8mm, inner sep=3pt},
  ext/.style={draw, align=center, fill=black!6, draw=black!45,
              minimum height=7mm, inner sep=3pt},
  clus/.style={->, thick, blue!70},
  capture/.style={->, very thick, blue!60},
  diss/.style={->, dashed, orange!85!black},
  rec/.style={<->, very thick, red!75},
  coal/.style={->, thick, orange!80!black},
  chain/.style={->, very thick, violet!70},
  snk/.style={->, thick, black!50},
  gbl/.style={->, very thick, green!45!black},
  src/.style={->, thick, green!55!black}]

  % ---- cascade source ---------------------------------------------------
  \node[ext] (G) at (-1.2,5.5) {cascade source $G$};

  % ---- the mobile block -------------------------------------------------
  \node[mob] (Cv)  at (-6.2,3.0) {$C_v$};
  \node[mob] (Ci)  at (-2.6,3.0) {$C_i$};
  \node[mob] (C2i) at (-0.2,3.0) {$C_{2i}$};
  \node[mob] (C3i) at ( 2.2,3.0) {$C_{3i}$};
  \node[draw, rounded corners=3pt, dashed, draw=blue!45, inner sep=6pt,
        fit=(Cv)(Ci)(C2i)(C3i), label={[font=\scriptsize, text=blue!60,
        anchor=south east]south east:{mobile species}}] (MOB) {};

  % ---- immobile families ------------------------------------------------
  \node[ipop] (ai) at (-3.4,0.3)
       {$(i,a_1),(i,a_2),(i,a_3)$\\[-1pt]
        \scriptsize prismatic $\langle a\rangle$, interstitial};
  \node[vpop] (av) at (-3.4,-2.2)
       {$(v,a_1),(v,a_2),(v,a_3)$\\[-1pt]
        \scriptsize prismatic $\langle a\rangle$, vacancy};
  \node[bpop] (c0) at (-5.6,-4.7) {$c_0$\\[-1pt]\scriptsize SFP};
  \node[bpop] (cf) at (-2.6,-4.7) {$c_f$\\[-1pt]\scriptsize faulted};
  \node[bpop] (cp) at ( 0.4,-4.7) {$c_p$\\[-1pt]\scriptsize perfect};

  % ---- external reservoirs ----------------------------------------------
  \node[ext] (NET) at (5.4,0.3)  {network\\[-1pt]\scriptsize $\rho_N$};
  \node[ext] (GB)  at (5.4,-3.0) {grain\\[-1pt]boundary};

  % ---- P4/P5 cascade sources --------------------------------------------
  \draw[src] (G.west) to[bend right=25]
       node[above left, pos=0.75, font=\scriptsize] {$G_m$ (P4)} (Cv.north);
  \draw[src] (G.south) -- node[right, pos=0.6, font=\scriptsize] {} (Ci.north);
  \draw[src] (G.south) -- (C2i.north);
  \draw[src] (G.south east) to[bend left=10] (C3i.north);
  \draw[src] (G.east) to[bend left=44]
       node[right, pos=0.35, font=\scriptsize] {$J_k$ (P5)} (ai.east);
  \draw[src] (G.east) to[bend left=52] (av.east);
  \draw[src, dashed] (G.east) to[bend left=58] (cf.east);

  % ---- P2 clustering, P3 dissociation ------------------------------------
  \draw[clus] (Ci.north east) to[bend left=30]
       node[above, font=\scriptsize] {$\beta^{i,i}$ (P2)} (C2i.north west);
  \draw[clus] (C2i.north east) to[bend left=30]
       node[above, font=\scriptsize] {$\beta^{i,2i}$} (C3i.north west);
  \draw[diss] (C2i.south west) to[bend left=30]
       node[below, font=\scriptsize] {$\alpha^{i,2i}$ (P3)} (Ci.south east);
  \draw[diss] (C3i.south west) to[bend left=30]
       node[below, font=\scriptsize] {$\alpha^{i,3i}$} (C2i.south east);

  % ---- P1 recombination and vacancy capture by clusters ------------------
  \draw[rec] (Cv) -- node[above, font=\scriptsize] {$R_{i,v}$ (P1)} (Ci);
  \draw[->, thick, red!60] (Cv.north east) to[bend left=22]
       node[above, pos=0.55, font=\scriptsize] {$R_{v,2i},R_{v,3i}$} (C3i.north west);

  % ---- P6 homogeneous loop nucleation -------------------------------------
  \draw[src, thick] (C3i.south) to[bend left=14]
       node[right, pos=0.55, font=\scriptsize]
       {$R_{i,3i}+R_{2i,2i}$ (P6)} (ai.north east);

  % ---- P7 capture at a loop or the SFP ------------------------------------
  \draw[capture] (MOB.south) to[bend right=6]
       node[left, pos=0.6, font=\scriptsize] {$\phi_{km}$ (P7)} (ai.north);
  \draw[capture] (MOB.south west) to[bend right=16] (av.north west);
  \draw[capture] (MOB.south west) to[bend right=34] (cf.north west);
  \draw[capture] (Cv.south) to[bend right=24] (c0.west);

  % ---- P8 peripheral emission ---------------------------------------------
  \draw[diss] (av.west) to[bend right=22]
       node[left, pos=0.5, font=\scriptsize] {$S_{vL}$ (P8)} (Cv.south west);
  \draw[diss] (cf.south) to[bend right=40] (Cv.south);

  % ---- P9 loop-loop coalescence -------------------------------------------
  \draw[coal] (ai.north west) to[out=150, in=210, looseness=5]
       node[left, font=\scriptsize] {$\varphi_{LL}$ (P9)} (ai.south west);
  \draw[coal] (cf.south west) to[out=-140, in=-40, looseness=5]
       node[below, font=\scriptsize] {$\varphi_{LL}$} (cf.south east);

  % ---- P10 loop-network, P11 mobile-network -------------------------------
  \draw[chain] (ai.east) -- node[above, font=\scriptsize]
       {$\varphi_{LN}$ (P10)} (NET.west);
  \draw[snk] (C3i.east) to[bend left=16]
       node[above, pos=0.7, font=\scriptsize] {$R_{m,s}$ (P11)} (NET.north);

  % ---- P12 basal chain ----------------------------------------------------
  \draw[chain] (c0) -- node[above, font=\scriptsize]
       {$\nu_{\mathrm{col}}$ (P12)} (cf);
  \draw[chain] (cf) -- node[above, font=\scriptsize]
       {$\nu_{\mathrm{uf}}$} (cp);

  % ---- P13 Dirichlet, P14 size-gated loop absorption ----------------------
  \draw[snk] (MOB.east) to[bend left=8]
       node[above, pos=0.6, font=\scriptsize] {Dirichlet (P13)} (GB.north);
  \draw[gbl] (av.east) to[bend right=6]
       node[above, pos=0.6, font=\scriptsize] {$\nu_{gb}$ (P14)} (GB.west);
  \draw[gbl] (cp.east) to[bend right=12] (GB.south);
\end{tikzpicture}
\caption{Reaction admissibility graph for the nine families of
\cref{eq:families}, with every edge class of \cref{tab:rag-processes}. The four
mobile species are boxed together because capture onto a loop sums over all of
them (\cref{eq:growth}), with a sign set by whether the species' polarity
matches the family's, so one bundled arrow per family is faithful where four
would be unreadable. Cascade production (green) feeds all four mobile species
and, as a separate partition of the same displacement rate, nucleates the
prismatic loops of both characters and the basal family directly; the dashed
source onto $c_f$ becomes a source onto $c_0$ instead when the basal chain is
active, which it is not in \cref{sec:results}. The interstitial prismatic
family has a second, homogeneous nucleation channel (P6) fed by
$2i+2i$ and $i+3i$ reactions among the mobile clusters, which is why the
mobile ladder and the loop populations are not separable. Dashed orange edges
run backwards from a population to the mobile pool: dissociation of a small
cluster (P3) and peripheral emission from a vacancy loop (P8). The three
prismatic variants of each character are drawn as one vertex for legibility;
they are separate populations that exchange nothing with one another and differ
only in habit normal, and the pyramid $c_0$ is not a loop --- it has no Burgers
vector, no perimeter and no coalescence channel, and enters only through a
capture cross-section.}
\label{fig:rag}
\end{figure}
